\par\medskip\noindent
2.8. \underline{Cambio de topos. Localización.}

Sea
\[
 f:\underline T'\longrightarrow\underline T
\]
un morfismo de topoi. El funtor de imagen inversa $f^*$ transforma una
biextensión de $(P,Q)$ por $G$ en una biextensión $f^*(E)$ de
$(f^*P,f^*Q)$ por $f^*G$, como sucede siempre para una estructura
expresable mediante límites proyectivos finitos y límites inductivos
arbitrarios. Para demostrar la compatibilidad de las dos estructuras
de extensión sobre $f^*(E)$ inducidas por las de $E$, se expresa la
compatibilidad mediante la conmutatividad de (2.1.1) en el caso
universal
\[
 S=P\times P\times Q\times Q,
\]
donde $p,p',q,q'$ son las cuatro proyecciones. Como era de esperar,
$E\mapsto f^*(E)$ da un funtor canónico
\[
 f^*:\operatorname{BIEXT}(\underline T)
 \longrightarrow
 \operatorname{BIEXT}(\underline T').
\]

En lo que sigue utilizaremos sin más mención el hecho de que toda
construcción hecha con biextensiones es compatible con las imágenes
inversas por morfismos de topoi. Esto se aplica, en particular, al
producto $EE'$ y al inverso $E^{-1}$ considerados en 2.5.

Se puede tomar como $\underline T'$ un \underline{topos inducido}
$\underline T_{/S'}$, y como $f$ el morfismo canónico, llamado
morfismo de localización (SGA 4 IV 5),
\[
 f=i_{S'}:\underline T_{/S'}\longrightarrow\underline T.
\]
Entonces $f^*(E)$ se denomina la \underline{biextensión inducida} por
$E$ sobre $S'$, y se denota también, como de costumbre, por $E|_{S'}$
o $E_{S'}$.

Como la restricción tiene las propiedades usuales de transitividad,
las categorías $\operatorname{BIEXT}(\underline T_{/S})$, para
$S\in\underline T$ variable, forman las fibras de una
\underline{categoría fibrada} sobre $\underline T$ (SGA 1 VI 8). De
hecho, esta categoría es incluso un \underline{stack} sobre
$\underline T$ en el sentido de Giraud [2]: los morfismos y los
objetos de sus fibras pueden pegarse. Por ejemplo, para dar una
biextensión sobre $S$ salvo isomorfismo canónico, basta dar una familia
de recubrimiento (es decir, una familia epimórfica)
$S_i\longrightarrow S$, una biextensión $E_i$ sobre cada $S_i$ (un
objeto de $\underline T_{/S_i}$) y, para cada par $(i,j)$, un
isomorfismo entre $E_i|_{S_j}$ y $E_j|_{S_i}$, sujeto a la condición de
cociclo usual.

Llamamos a este el \underline{stack de biextensiones sobre
$\underline T$}. Sustituyendo $\underline T$ en (2.3.3), y por tanto
también en (2.3.4) y (2.3.5), por $\underline T_{/S}$ para $S$
variable, se obtienen morfismos de stacks sobre $\underline T$. Lo
mismo vale para las operaciones $E.E'$ y $E^{-1}$ de (2.5.1) en las
categorías $\operatorname{BIEXT}(P_S,Q_S;G_S)$ para $S$ variable,
donde $P,Q,G$ son grupos conmutativos fijos sobre $\underline T$.
